% Generated by scripts/gen_blueprint.py from blueprint/nodes/04-gelfond-schneider.toml. Do not edit.

\chapter{Gelfond–Schneider}
\label{chap:04-gelfond-schneider}

The Gelfond--Schneider theorem in logarithmic form, by Gel'fond's method, and the transcendence of $e^{\pi}$. The Lean proof restructures the formalisation of M.~Karatarakis and F.~Wiedijk \cite{KW26} (Apache 2.0), which the library first carried as a single module of 5,388 lines. It is now a tree of eleven results, 1,409 lines in all: two tools of Chapter~\ref{chap:02-tools}, Lemmas~\ref{lem:liouville_house} and~\ref{lem:expSum_first_nonvanishing}; the eight steps of Gel'fond's method below; and Theorem~\ref{thm:gelfond_schneider}. Lemma~\ref{lem:deriv_upper} goes through the grid estimate of Lemma~\ref{lem:expPoly_grid_estimate}, which the four exponentials development shares.

Unless said otherwise, $K$ is a number field, $\sigma\colon K\to\C$ a field embedding and $\alpha',\beta',\gamma'\in K$; sums over $a,b$ run over $1\le a,b\le q$.

% Prove2Me: GelfondSchneider.common_field -- https://prove2.me/theorems/225f1e21-3d8c-442b-a6c8-153398b69c97
\begin{lemma}[One number field]
\label{lem:common_field}
\lean{GelfondSchneider.common_field}
\leanok
Let $l,\beta\in\C$ with $\beta$, $e^{l}$ and $e^{\beta l}$ algebraic. There are a number field $K$, a field embedding $\sigma\colon K\to\C$ and $\alpha',\beta',\gamma'\in K$ with $\sigma(\alpha')=e^{l}$, $\sigma(\beta')=\beta$ and $\sigma(\gamma')=e^{\beta l}$.

{\small\emph{Source:} A step of Gel'fond's proof \cite{Gel34}.}
\end{lemma}

\begin{proof}
\leanok
Take $K=\Q(e^{l},\beta,e^{\beta l})$ and its inclusion in $\C$.
\end{proof}

% Prove2Me: GelfondSchneider.system_entry_house_le -- https://prove2.me/theorems/cbe83e8d-f145-4454-a29e-dcf7bd32ac01
\begin{lemma}[Entries of the linear system]
\label{lem:system_entry_house_le}
\lean{GelfondSchneider.system_entry_house_le}
\leanok
Let $m\ge1$. There is $C\ge1$ such that for all $n\ge1$ and $q$ with $q^{2}=2mn$, all $1\le a,b\le q$, $1\le j\le m$ and $k<n$,
\[ \houseof{(a+b\beta')^{k}\alpha'^{\,aj}\gamma'^{\,bj}}\le C^{n}n^{(n-1)/2}. \]

{\small\emph{Source:} A step of Gel'fond's proof \cite{Gel34}.}
\end{lemma}

\begin{proof}
\leanok
\end{proof}

% Prove2Me: GelfondSchneider.aux_coeffs -- https://prove2.me/theorems/58990e61-4127-41a3-b61f-f9254333c532
\begin{lemma}[Gel'fond's auxiliary function]
\label{lem:aux_coeffs}
\lean{GelfondSchneider.aux_coeffs}
\leanok
Let $\alpha',\gamma'\neq0$ and $m\ge1$. There is $C\ge1$ such that for all $n\ge1$ and $q$ with $q^{2}=2mn$ there are algebraic integers $\eta_{ab}\in\mathcal{O}_K$, not all zero, with $\houseof{\eta_{ab}}\le C^{n}n^{(n+1)/2}$ and
\[ \sum_{a,b}\eta_{ab}(a+b\beta')^{k}\alpha'^{\,aj}\gamma'^{\,bj}=0\qquad(1\le j\le m,\ 0\le k<n). \]

{\small\emph{Source:} A step of Gel'fond's proof \cite{Gel34}.}
\end{lemma}

\begin{proof}
\leanok
\uses{lem:system_entry_house_le}
There are $q^{2}=2mn$ unknowns and $mn$ equations, so Mathlib's Siegel lemma over $\mathcal{O}_K$ applies with exponent one once denominators are cleared; Lemma~\ref{lem:system_entry_house_le} bounds the entries.
\end{proof}

% Prove2Me: GelfondSchneider.deriv_identity -- https://prove2.me/theorems/c21422a2-1acb-478f-8fd1-ae8ea5b082f4
\begin{lemma}[Derivatives at the integers]
\label{lem:deriv_identity}
\lean{GelfondSchneider.deriv_identity}
\leanok
Let $K$ be any field, $\sigma\colon K\to\C$ a ring homomorphism, and $l,\beta\in\C$ with $\sigma(\beta')=\beta$, $\sigma(\alpha')=e^{l}$ and $\sigma(\gamma')=e^{\beta l}$. For $\eta_{ab}\in K$ put $E(z)=\sum_{a,b}\sigma(\eta_{ab})e^{(a+b\beta)lz}$. Then for all $k,j\in\N$
\[ E^{(k)}(j)=l^{k}\,\sigma\Bigl(\sum_{a,b}\eta_{ab}(a+b\beta')^{k}\alpha'^{\,aj}\gamma'^{\,bj}\Bigr). \]

{\small\emph{Source:} A step of Gel'fond's proof \cite{Gel34}.}
\end{lemma}

\begin{proof}
\leanok
Differentiate term by term: $e^{(a+b\beta)lj}=\sigma(\alpha'^{\,aj}\gamma'^{\,bj})$.
\end{proof}

% Prove2Me: GelfondSchneider.rho_denominator -- https://prove2.me/theorems/cd68ef56-23e3-4b24-9858-0f3a3818c994
\begin{lemma}[A common denominator]
\label{lem:rho_denominator}
\lean{GelfondSchneider.rho_denominator}
\leanok
Let $m\in\N$. There is an integer $D\neq0$ such that for all $n,q,r,l_0$ with $n\ge1$, $q^{2}=2mn$, $n\le r$ and $1\le l_0\le m$, and all $\eta_{ab}\in\mathcal{O}_K$, the number
\[ D^{r}\sum_{a,b}\eta_{ab}(a+b\beta')^{r}\alpha'^{\,al_0}\gamma'^{\,bl_0} \]
is an algebraic integer.

{\small\emph{Source:} A step of Gel'fond's proof \cite{Gel34}.}
\end{lemma}

\begin{proof}
\leanok
The powers of $\alpha'$, $\beta'$ and $\gamma'$ that occur have exponents at most linear in $r$, since $q\le q^{2}=2mn\le2mr$.
\end{proof}

% Prove2Me: GelfondSchneider.rho_house_le -- https://prove2.me/theorems/62c29048-d0cf-49fc-8309-6427246ebbcf
\begin{lemma}[House of the first non-zero derivative]
\label{lem:rho_house_le}
\lean{GelfondSchneider.rho_house_le}
\leanok
Let $m\ge1$ and $C_0\ge1$. There is $C\ge1$ such that whenever $n\ge1$, $q^{2}=2mn$, $n\le r$, $1\le l_0\le m$, and $\eta_{ab}\in\mathcal{O}_K$ satisfy $\houseof{\eta_{ab}}\le C_0^{n}n^{(n+1)/2}$,
\[ \houseof{\sum_{a,b}\eta_{ab}(a+b\beta')^{r}\alpha'^{\,al_0}\gamma'^{\,bl_0}}\le C^{r}r^{\,r+3/2}. \]

{\small\emph{Source:} A step of Gel'fond's proof \cite{Gel34}.}
\end{lemma}

\begin{proof}
\leanok
\end{proof}

% Prove2Me: GelfondSchneider.deriv_upper -- https://prove2.me/theorems/d74690db-d2c0-4b78-81b2-800698673390
\begin{lemma}[Upper bound for the first non-zero derivative]
\label{lem:deriv_upper}
\lean{GelfondSchneider.deriv_upper}
\leanok
Let $l,\beta\in\C$, $m\ge1$ and $C_0\ge1$. There is $C\ge1$ with the following property. Let $n\ge1$, $q^{2}=2mn$, $n\le r$ and $1\le l_0\le m$, and let $E(z)=\sum_{a,b}c_{ab}e^{(a+b\beta)lz}$ with $c_{ab}\in\C$, $|c_{ab}|\le C_0^{n}n^{(n+1)/2}$. If $E^{(k)}(j)=0$ for all $1\le j\le m$ and $k<r$, then
\[ |E^{(r)}(l_0)|\le C^{r}r^{\,r(3-m)/2+3/2}. \]

{\small\emph{Source:} A step of Gel'fond's proof \cite{Gel34}.}
\end{lemma}

\begin{proof}
\leanok
\uses{lem:expPoly_grid_estimate}
Lemma~\ref{lem:expPoly_grid_estimate} for the one-dimensional grid of the points $1,\dots,m$: the zeros of order $r$ at all of them, $l_0$ included, give the factor $r^{-rm/2}$.
\end{proof}

% Prove2Me: GelfondSchneider.main_estimate -- https://prove2.me/theorems/afd443ea-1ee0-45e7-bf29-99baf10f15ef
\begin{lemma}[Gel'fond's main estimate]
\label{lem:main_estimate}
\lean{GelfondSchneider.main_estimate}
\leanok
Let $K$ have degree $h$, let $l\in\C$, $l\neq0$, and let $\beta\in\C\setminus\Q$, with $\sigma(\alpha')=e^{l}$, $\sigma(\beta')=\beta$ and $\sigma(\gamma')=e^{\beta l}$. There is $C\ge1$ such that for every $N\in\N$ there is an integer $r\ge\max(N,1)$ with
\[ r^{(r-3h)/2}\le C^{r}. \]

{\small\emph{Source:} A step of Gel'fond's proof \cite{Gel34}.}
\end{lemma}

\begin{proof}
\leanok
\uses{lem:aux_coeffs,lem:deriv_identity,lem:expSum_first_nonvanishing,lem:rho_denominator,lem:rho_house_le,lem:deriv_upper,lem:liouville_house}
Take $m=2h+2$ and the auxiliary function of Lemma~\ref{lem:aux_coeffs}; its frequencies are distinct because $\beta\notin\Q$. Lemma~\ref{lem:expSum_first_nonvanishing} gives a first non-zero derivative $E^{(r)}(l_0)$, which by Lemma~\ref{lem:deriv_identity} is $l^{r}$ times the image of an element of $K$. Lemmas~\ref{lem:rho_denominator}, \ref{lem:rho_house_le} and~\ref{lem:liouville_house} bound it from below, and Lemma~\ref{lem:deriv_upper} from above.
\end{proof}

% Prove2Me: Schanuel.gelfond_schneider -- https://prove2.me/theorems/2e4b9062-99bc-4797-b90c-279b63f084aa
\begin{theorem}[Gelfond–Schneider]
\label{thm:gelfond_schneider}
\lean{GelfondSchneider.gelfond_schneider}
\leanok
Let $l\in\C$, $l\neq0$, with $e^{l}$ algebraic, and let $\beta$ be an algebraic number that is not rational. Then $e^{\beta l}$ is transcendental.

{\small\emph{Source:} Gel'fond \cite{Gel34} and Schneider \cite{Sch34}, Hilbert's seventh problem; see \cite[Theorem 2.1]{Bak75}.}
\end{theorem}

\begin{proof}
\leanok
\uses{lem:common_field,lem:main_estimate}
If $e^{\beta l}$ were algebraic, Lemma~\ref{lem:common_field} would put $e^{l}$, $\beta$ and $e^{\beta l}$ in one number field, and Lemma~\ref{lem:main_estimate} would give $r^{(r-3h)/2}\le C^{r}$ for arbitrarily large $r$, which is false.
\end{proof}

% Prove2Me: e_pi_transcendence -- https://prove2.me/theorems/e958f5e6-60e1-4bf4-9dd7-6e6864044301
% Statement written from the Lean source: the page text is a question with commentary, not a statement
\begin{theorem}[Transcendence of $e^{\pi}$]
\label{thm:e_pi_transcendence}
\lean{Diaz.e_pi_transcendence}
\leanok
The number $e^{\pi}$ is transcendental.

{\small\emph{Source:} Gel'fond (1929).}
\end{theorem}

\begin{proof}
\leanok
\uses{thm:gelfond_schneider}
Theorem~\ref{thm:gelfond_schneider} with $l=i\pi$ and $\beta=-i$: $e^{i\pi}=-1$ is algebraic, and $e^{\beta l}=e^{\pi}$.
\end{proof}
